\subsection{Usability — System Usability Scale}
\label{sec:eval-sus}

The System Usability Scale captures how usable each interface felt
to participants immediately after they finished interacting with it.
It consists of ten Likert items on a five-point scale, alternating in
polarity: odd-numbered items are phrased positively (higher response
equals more positive evaluation), and even-numbered items are phrased
negatively (lower response equals more positive evaluation). This
alternating keying is why the per-item table that follows must be
read item by item rather than averaged at a glance. The full item
text is reproduced in Appendix~\ref{app:sus-rtlx-items}. The
aggregated ``All'' column in the tables below pools the three
per-mode samples into a single per-item descriptive across all $207$
ratings.

The ten SUS items capture distinct facets of usability — frequency of
intended use, perceived complexity, ease of use, and so on — and the
per-mode means therefore have to be read item by item before the
aggregated test is reported. Table~\ref{tab:sus-per-item} reports
the mean response per item per mode and the aggregated mean across
all three modes. The aggregated column is the unweighted mean of the
three per-mode means.

\begin{table}[htbp]
\centering
\caption{System Usability Scale, mean response per item by mode ($n = 69$ per
mode). Items are reported on the original five-point Likert scale; the
``All'' column averages the three per-mode means. Odd items: higher is better;
even items: lower is better.}
\label{tab:sus-per-item}
\small
\begin{tabularx}{\linewidth}{@{}c X r r r r@{}}
\toprule
\# & Item & Form & Ask & Chat & All \\
\midrule
1 & I think that I would like to use this system frequently. & 3.52 & 3.22 & 3.22 & 3.32 \\
2 & I found the system unnecessarily complex. & 1.64 & 2.25 & 2.12 & 2.00 \\
3 & I thought the system was easy to use. & 4.38 & 4.01 & 3.99 & 4.13 \\
4 & I think that I would need the support of a technical person to be able to use this system. & 1.52 & 1.49 & 1.55 & 1.52 \\
5 & I found the various functions in this system were well integrated. & 4.00 & 3.67 & 3.64 & 3.77 \\
6 & I thought there was too much inconsistency in this system. & 2.16 & 2.04 & 2.01 & 2.07 \\
7 & I would imagine that most people would learn to use this system very quickly. & 4.30 & 4.00 & 4.03 & 4.11 \\
8 & I found the system very cumbersome to use. & 1.65 & 2.16 & 2.29 & 2.03 \\
9 & I felt very confident using the system. & 3.83 & 3.84 & 3.90 & 3.86 \\
10 & I needed to learn a lot of things before I could get going with this system. & 1.61 & 1.61 & 1.62 & 1.61 \\
\bottomrule
\end{tabularx}
\end{table}

Form is the mode participants would most willingly keep using, with
Ask and Chat trailing by a small but consistent margin. For
\emph{``I think that I would like to use this system frequently''}
(item~1), Form yields the highest mean ($3.52$), followed by Ask
and Chat at the same level ($3.22$ each); aggregated mean $3.32$.
Frequent-use intent decreases by approximately $0.30$~points from
Form to the LLM-driven modes.

Form also feels the least complex of the three, with Ask perceived
as marginally more complex than Chat. For \emph{``I found the
system unnecessarily complex''} (item~2, lower is better), Form is
rated lowest ($1.64$), Chat second ($2.12$), and Ask highest at
$2.25$; aggregated mean $2.00$. Perceived complexity is
approximately $0.5$~points higher in the LLM-driven modes than in
Form.

Form is also experienced as the easiest mode to use, with Ask and
Chat almost indistinguishable behind it. For \emph{``I thought the
system was easy to use''} (item~3), Form is rated highest ($4.38$),
with Ask and Chat closely matched at $4.01$ and $3.99$; aggregated
mean $4.13$. Ease-of-use ratings decrease by roughly $0.4$~points
from Form to the LLM-driven modes.

None of the three modes makes participants feel they need technical
help to operate it — the Form-first ordering levels out here. For
\emph{``I think that I would need the support of a technical person
to be able to use this system''} (item~4, lower is better), the
three modes are essentially indistinguishable (Form $1.52$, Ask
$1.49$, Chat $1.55$), with an aggregated mean of $1.52$. None of
the three modes is perceived to require technical support.

Form feels like the most coherent mode, with the LLM-driven modes
scoring lower on perceived integration. For \emph{``I found the
various functions in this system were well integrated''} (item~5),
Form scores highest ($4.00$), Ask and Chat lower ($3.67$ and
$3.64$); aggregated mean $3.77$. The integration rating drops by
approximately $0.35$~points from Form to the LLM-driven modes.

No mode is perceived as meaningfully more inconsistent than the
others, mildly inverting the Form-first pattern. For \emph{``I
thought there was too much inconsistency in this system''} (item~6,
lower is better), the three modes are similar (Form $2.16$, Ask
$2.04$, Chat $2.01$), aggregated mean $2.07$. Form is rated
marginally higher on this negative item than the two LLM-driven
modes.

Form is judged the most learnable of the three, with Ask and Chat
trailing again. For \emph{``I would imagine that most people would
learn to use this system very quickly''} (item~7), Form scores
$4.30$, Ask $4.00$, Chat $4.03$; aggregated mean $4.11$. Perceived
learnability decreases by approximately $0.30$~points from Form to
the LLM-driven modes.

The biggest experiential gap of the ten items lands here:
participants find Form noticeably less cumbersome than the
LLM-driven modes, with Chat the most cumbersome of the three. For
\emph{``I found the system very cumbersome to use''} (item~8,
lower is better), the largest spread is observed: Form $1.65$, Ask
$2.16$, Chat $2.29$; aggregated mean $2.03$. Cumbersomeness is
rated approximately $0.5$--$0.65$~points higher in the LLM-driven
modes than in Form.

Participants feel about equally confident operating any of the
three modes, with a small uptick toward Chat — the only item where
the LLM-driven modes nudge ahead of Form. For \emph{``I felt very
confident using the system''} (item~9), all three modes are rated
similarly and slightly higher in the LLM-driven modes (Form $3.83$,
Ask $3.84$, Chat $3.90$); aggregated mean $3.86$. Within-system
confidence is broadly comparable across modes, with a small upward
trend from Form to Chat.

All three modes feel equally approachable to a first-time user — no
prior learning is perceived as required for any of them. For
\emph{``I needed to learn a lot of things before I could get going
with this system''} (item~10, lower is better), the three modes are
virtually identical (Form $1.61$, Ask $1.61$, Chat $1.62$),
aggregated mean $1.61$. The amount of prior learning required is
perceived as low and equal across all modes.

Table~\ref{tab:sus-per-item-distribution} reports the full
descriptive distribution per item per mode (mean, standard deviation,
median, Q1, Q3). The minimum and maximum are $1$ and $5$ for every
item~$\times$~mode combination and are omitted for compactness.

\begin{table}[htbp]
\centering
\caption{System Usability Scale, full per-item descriptive
distribution by mode ($n = 69$ per mode). Source:
\texttt{sus\_per\_item\_summary.csv}.}
\label{tab:sus-per-item-distribution}
\small
\begin{tabular}{@{}c l r r r r r@{}}
\toprule
\# & Mode & Mean & SD & Median & Q1 & Q3 \\
\midrule
\multirow{4}{*}{1}  & All  & 3.32 & 1.14 & 3 & 3 & 4 \\
                    & Form & 3.52 & 1.09 & 4 & 3 & 4 \\
                    & Ask  & 3.22 & 1.08 & 3 & 3 & 4 \\
                    & Chat & 3.22 & 1.24 & 3 & 2 & 4 \\
\midrule
\multirow{4}{*}{2}  & All  & 2.00 & 1.12 & 2 & 1 & 3 \\
                    & Form & 1.64 & 0.92 & 1 & 1 & 2 \\
                    & Ask  & 2.25 & 1.22 & 2 & 1 & 3 \\
                    & Chat & 2.12 & 1.12 & 2 & 1 & 3 \\
\midrule
\multirow{4}{*}{3}  & All  & 4.13 & 0.97 & 4 & 4 & 5 \\
                    & Form & 4.38 & 0.91 & 5 & 4 & 5 \\
                    & Ask  & 4.01 & 0.81 & 4 & 4 & 5 \\
                    & Chat & 3.99 & 1.13 & 4 & 3 & 5 \\
\midrule
\multirow{4}{*}{4}  & All  & 1.52 & 0.89 & 1 & 1 & 2 \\
                    & Form & 1.52 & 0.85 & 1 & 1 & 2 \\
                    & Ask  & 1.49 & 0.80 & 1 & 1 & 2 \\
                    & Chat & 1.55 & 1.01 & 1 & 1 & 2 \\
\midrule
\multirow{4}{*}{5}  & All  & 3.77 & 1.05 & 4 & 3 & 5 \\
                    & Form & 4.00 & 0.92 & 4 & 4 & 5 \\
                    & Ask  & 3.67 & 1.02 & 4 & 3 & 4 \\
                    & Chat & 3.64 & 1.16 & 4 & 3 & 5 \\
\midrule
\multirow{4}{*}{6}  & All  & 2.07 & 1.00 & 2 & 1 & 3 \\
                    & Form & 2.16 & 0.96 & 2 & 1 & 3 \\
                    & Ask  & 2.04 & 1.06 & 2 & 1 & 3 \\
                    & Chat & 2.01 & 0.98 & 2 & 1 & 3 \\
\midrule
\multirow{4}{*}{7}  & All  & 4.11 & 0.98 & 4 & 4 & 5 \\
                    & Form & 4.30 & 0.96 & 5 & 4 & 5 \\
                    & Ask  & 4.00 & 0.95 & 4 & 4 & 5 \\
                    & Chat & 4.03 & 1.01 & 4 & 4 & 5 \\
\midrule
\multirow{4}{*}{8}  & All  & 2.03 & 1.13 & 2 & 1 & 3 \\
                    & Form & 1.65 & 0.94 & 1 & 1 & 2 \\
                    & Ask  & 2.16 & 1.12 & 2 & 1 & 3 \\
                    & Chat & 2.29 & 1.23 & 2 & 1 & 3 \\
\midrule
\multirow{4}{*}{9}  & All  & 3.86 & 1.04 & 4 & 3 & 5 \\
                    & Form & 3.83 & 1.04 & 4 & 3 & 5 \\
                    & Ask  & 3.84 & 0.93 & 4 & 3 & 4 \\
                    & Chat & 3.90 & 1.15 & 4 & 4 & 5 \\
\midrule
\multirow{4}{*}{10} & All  & 1.61 & 0.90 & 1 & 1 & 2 \\
                    & Form & 1.61 & 0.86 & 1 & 1 & 2 \\
                    & Ask  & 1.61 & 0.81 & 1 & 1 & 2 \\
                    & Chat & 1.62 & 1.03 & 1 & 1 & 2 \\
\bottomrule
\end{tabular}
\end{table}

The per-item distributions are visualised in Figure~\ref{fig:sus-subquestions}.

\begin{sidewaysfigure}
\centering
\includegraphics[width=\linewidth]{sus_subquestions_boxplot.png}
\caption{Box plots of the ten SUS items by mode ($n = 69$ per mode). Boxes show
the interquartile range; the central line marks the median, and the diamond
marks the mean.}
\label{fig:sus-subquestions}
\end{sidewaysfigure}

A Friedman test on the per-mode aggregate yields
$\chi^2(2) = 5.51$, $p = 0.064$, $W = 0.040$, which does not reach
the $\alpha = 0.05$ threshold. Pairwise Wilcoxon signed-rank post-hoc
tests with Bonferroni correction ($\alpha = 0.0167$) indicate a
significant Form-vs-Ask difference ($W = 578.5$, $p = 0.013$,
rank-biserial $r = +0.368$, medium effect), a non-significant
Form-vs-Chat difference ($p = 0.018$, just above the Bonferroni
threshold), and no Ask-vs-Chat difference ($p = 0.666$). Full test
statistics are reproduced in Appendix~\ref{app:inferential}.

Taken together, the SUS data establish a Form-first ordering on the
items most directly tied to perceived ease of use, complexity, and
cumbersomeness, while items relating to within-system confidence,
required prior learning, and need for technical support flatten
across modes. Whether this same ordering recurs on perceived
workload is examined next in \S\ref{sec:eval-rtlx}.
